\chapter{Build, Continuous Integration, and Verification Gates}
\label{ch:build}

A safety-boundary architecture is only as trustworthy as the gate that
enforces it. This chapter describes the build system and the verification
gates that run the pipeline end to end.

\section{Build Targets}
\label{sec:build-targets}

The repository is multi-language and multi-toolchain:
\begin{itemize}
  \item \textbf{Rust/Cargo} (\texttt{Cargo.toml}) for the orchestration
        surface.
  \item \textbf{Bazel} (\texttt{WORKSPACE.bazel}) for the LLVM/C{--} backend.
  \item \textbf{Nix} (\texttt{flake.nix}) for reproducible dev environments.
  \item \textbf{Lean} (\texttt{lakefile.lean}, \texttt{lean-toolchain}) for
        the proof layer.
  \item \textbf{Shell/CI} for the gates.
\end{itemize}

\section{The End-to-End Scaffold}
\label{sec:build-e2e}

\texttt{ci/run\_e2e.sh} sketches the full pipeline as a linear script:
\begin{lstlisting}[language=BashX,caption={End-to-end scaffold (run\_e2e.sh).},label=lst:build-e2e]
# 1. Parse APL -> S-expression
# 2. Normalize S-expression
# 3. Liquid Haskell refinement + Lean obligation gen
# 4. Lean proves obligations (or fails)
# 5. Closed-form transformation
# 6. TeX doc generation
# 7. LLVM backend -> object file
# 8. Link with verified FFI shim
echo "Pipeline scaffold verified for $APL_FILE"
\end{lstlisting}
The steps are currently commented (the individual layer tools exist; the
orchestration glue is scaffolding), but the ordering encodes the contract
sequence of Chapters~\ref{ch:apl}--\ref{ch:completeness} exactly.

\section{The Verification Gate}
\label{sec:build-gate}

\texttt{mathlib5-ffi-bridge/scripts/verify.sh} is the executable safety
boundary. It runs five tests and fails the build if any do:
\begin{enumerate}
  \item C bridge compiles (\texttt{clang -shared -fPIC}).
  \item Lean builds (\texttt{lake build}).
  \item FFI round-trip passes (greps \texttt{All tests passed}).
  \item Theorem calculator builds (\texttt{lake build theorem\_calc}).
  \item Receipt is generated and sealed.
\end{enumerate}
The gate tallies \texttt{PASS} and \texttt{FAIL}; its final lines print
\texttt{Status: VERIFIED} or \texttt{FAILED} and it exits with the failure
count, so CI treats a non-zero exit as a broken boundary.

\section{Receipts in CI}
\label{sec:build-receipt}

After the gate, \texttt{receipt.sh} (\cref{sec:worm-receipt}) hashes sources
and binary, writes \texttt{manifest.json} and \texttt{seal.sha256}, and stamps
the git SHA and Lean toolchain. The receipt is the artifact that Layer~9
appends to the WORM ledger, closing the provenance loop for the build itself.

\section{Reproducibility}
\label{sec:build-repro}

Three mechanisms protect reproducibility:
\begin{itemize}
  \item \textbf{Pinned toolchain} --- \texttt{lean-toolchain} fixes the Lean
        version; the receipt records it.
  \item \textbf{Nix flake} --- \texttt{flake.nix} locks the entire dependency
        graph for dev and CI.
  \item \textbf{Content-addressed inputs} --- every layer consumes hash-stable
        artifacts, so re-running produces identical receipts.
\end{itemize}

\section{The Gate as a Boundary}
\label{sec:build-asboundary}

Crucially, the verification gate is itself a safety boundary with an inbound
contract (``source tree'') and an outbound contract (``a sealed, verified
receipt or a hard failure''). There is no ``mostly passing'' state: a single
failed test fails the build. This is what makes the pipeline's guarantees
operational rather than aspirational.
